% This file is a chapter fragment.  It is included by version/detailed.tex.

\chapter{Zeta linear algebra}\label{chap:zla}

This chapter is the purely algebraic endgame of the Weil conjectures: what one
can deduce about a power series
\[
  Z(T)=\operatorname{Exp}\left(\sum_{n\geq1}N_n\frac{T^n}{n}\right)
\]
from the sole knowledge that its counting numbers $N_n$ are the alternating
sums of the traces of the iterates of an endomorphism on a finite family of
finite-dimensional vector spaces.  No geometry and no cohomology enter here; the
input is a \emph{trace datum} (\Cref{def:zla-trace-datum}), which is what the
Grothendieck--Lefschetz trace formula will eventually supply.

The two statements this material serves are
\begin{itemize}
\item \texttt{AlgebraicGeometry.Scheme.zetaSeries\_rational} of
      \texttt{WeilConjectures/Statement.lean}: the zeta series of a scheme of
      finite type over a finite field is $P/Q$ with $P,Q\in\mathbf Z[T]$ and
      $Q(0)=1$.  The chain that produces it is
      \Cref{prop:zla-cohomological-factorization} (rationality after scalar
      extension), \Cref{thm:zla-hankel-descent} together with
      \Cref{cor:zla-rationality-counting-field} (descent of rationality to
      $\mathbf Q$), and \Cref{thm:zla-integral-fatou} (integrality of the
      reduced fraction).
\item \texttt{AlgebraicGeometry.IsWeilFactorization} of
      \texttt{Comparator/Weil/Definitions.lean}: the factorization
      $Z\cdot\prod_{i\text{ even}}P_i=\prod_{i\text{ odd}}P_i$ with
      $P_i\in\mathbf Z[T]$, $P_i(0)=1$.  Its factorization identity comes from
      \Cref{prop:zla-cohomological-factorization}, and the integrality of the
      individual $P_i$ from
      \Cref{prop:zla-parity-identification},
      \Cref{prop:zla-unique-recovery} and
      \Cref{prop:zla-factor-integrality}, which separate the factors by
      weight.  The last section prepares the functional equation by computing
      the reciprocal characteristic polynomial of the partner of a Frobenius
      under a perfect similitude pairing
      (\Cref{thm:zla-degreewise-reciprocal}).
\end{itemize}

The source is the sketch article
\texttt{blueprint/src/version/sketch/49\_zeta\_linear\_algebra.tex}, and only
the statements that are to be formalised in this subtree are reproduced here:
its sections on formal power series at the origin, locally finite products,
periodic permutations, closed points and the Hasse--Weil series, compatible
realizations, the middle pairing and its determinant sign, bounded graded
duality systems, the raw and normalized functional equations,
realization-independent constants and the final arithmetic theorem are
deliberately omitted, as are the derived-category formulation of a trace datum
and the statements about perfect complexes that go with it.

Throughout, $R[[T]]$ is the ring of formal power series over a commutative
ring $R$, identities in it are coefficientwise, and for a commutative
$\mathbf Q$-algebra $R$ we write
\[
  \operatorname{Exp}(U)=\sum_{r\geq0}\frac{U^r}{r!}\in1+TR[[T]],
  \qquad
  \operatorname{Log}(1+V)=\sum_{r\geq1}(-1)^{r+1}\frac{V^r}{r}\in TR[[T]]
\]
for $U,V\in TR[[T]]$; both sums are coefficientwise finite.

\section{Determinants and traces in finite dimension}\label{sec:zla-det-trace}

Let $E$ be a field of characteristic zero.

\begin{definition}[Reciprocal characteristic polynomial]
\label{def:zla-recip-charpoly}
For a finite-dimensional $E$-vector space $V$ and
$u\in\operatorname{End}_E(V)$, define
\[
  P_{V,u}(T)=\det(1-Tu\mid V)\in E[T].
\]
After scalar extension to an algebraic closure $\Omega/E$, write
\[
  P_{V,u}(T)=\prod_{j=1}^{\dim_EV}(1-\alpha_jT).
\]
The multiset $\{\alpha_j\}$, with algebraic multiplicities, is the multiset of
\emph{reciprocal roots} of $P_{V,u}$; it is the eigenvalue multiset of $u$.
If $u$ is invertible then every $\alpha_j\neq0$, $\deg P_{V,u}=\dim_EV$, and
the leading coefficient is
$[T^{\dim V}]P_{V,u}=(-1)^{\dim V}\det(u)$.
\end{definition}

\begin{lemma}[Finite-dimensional trace--determinant identity]
\label{lem:zla-trace-det}
\uses{def:zla-recip-charpoly}
Let $V$ be a finite-dimensional $E$-vector space and let
$u\in\operatorname{End}_E(V)$.  Then
\[
  P_{V,u}(T)^{-1}
  =\operatorname{Exp}\left(\sum_{n\geq1}\Tr(u^n\mid V)\frac{T^n}{n}\right)
  \qquad\text{in }E[[T]],
\]
the inverse on the left being taken in $E[[T]]$, which is legitimate because
$P_{V,u}(0)=1$.  Equivalently,
\[
  T\frac{d}{dT}\operatorname{Log}P_{V,u}(T)^{-1}
  =\sum_{n\geq1}\Tr(u^n\mid V)T^n.
\]
No semisimplicity assumption is imposed on $u$.
\end{lemma}
\begin{proof}
Differentiating a determinant gives the matrix identity
$\frac{d}{dT}\det(1-Tu)=-\det(1-Tu)\Tr\!\left((1-Tu)^{-1}u\right)$ in
$E[[T]]$, where the formal resolvent is
$(1-Tu)^{-1}=\sum_{n\geq0}T^nu^n$; taking the trace coefficientwise turns the
right-hand side into $-\det(1-Tu)\sum_{n\geq1}\Tr(u^n)T^{n-1}$.  Both sides of
the asserted identity therefore have constant coefficient $1$ and the same
logarithmic derivative, and over a $\mathbf Q$-algebra a power series is
determined by its constant coefficient and its derivative.
\end{proof}

\section{The trace datum and cohomological factorization}
\label{sec:zla-trace-datum}

\begin{definition}[Finite cohomological trace datum]
\label{def:zla-trace-datum}
\uses{def:zla-recip-charpoly}
Let $E$ be a field of characteristic zero.  A \emph{finite cohomological trace
datum} over $E$ consists of:
\begin{enumerate}[label=\textup{(\roman*)}]
\item an integer $r\geq0$ and, for each $0\leq i\leq r$, a finite-dimensional
      $E$-vector space $V_i$ together with an endomorphism
      $F_i\in\operatorname{End}_E(V_i)$;
\item integers $N_n\in\mathbf Z$ for every $n\geq1$;
\item the equalities
\[
  N_n=\sum_{i=0}^r(-1)^i\Tr(F_i^n\mid V_i)
  \qquad\text{in }E,\qquad(n\geq1),
\]
the left-hand side being the image of $N_n$ under
$\mathbf Z\to\mathbf Q\to E$.
\end{enumerate}
Put $P_i(T)=\det(1-TF_i\mid V_i)$.  In the intended application $V_i$ is the
$i$-th cohomology group of a perfect complex and $F_i$ is induced by a
Frobenius automorphism, but nothing below uses a complex: only the finite
family $(V_i,F_i)_{0\leq i\leq r}$ and the trace equalities enter.
\end{definition}

\begin{lemma}[Counting logarithm equals trace logarithm]
\label{lem:zla-counting-log}
\uses{def:zla-trace-datum}
For a finite cohomological trace datum,
\[
  \sum_{n\geq1}N_n\frac{T^n}{n}
  =\sum_{n\geq1}\left(\sum_{i=0}^r(-1)^i\Tr(F_i^n\mid V_i)\right)\frac{T^n}{n}
  \qquad\text{in }TE[[T]].
\]
\end{lemma}
\begin{proof}
The two series are compared coefficientwise; in degree $n$ the assertion is
the image in $E$ of the equality (iii) of \Cref{def:zla-trace-datum} divided by
the invertible integer $n$.  The left-hand side is the coefficientwise scalar
extension along $\mathbf Q\to E$ of a series with rational coefficients.
\end{proof}

\begin{proposition}[Cohomological determinant factorization]
\label{prop:zla-cohomological-factorization}
\uses{def:zla-trace-datum, lem:zla-counting-log, lem:zla-trace-det}
For a finite cohomological trace datum,
\[
  \operatorname{Exp}\left(\sum_{n\geq1}N_n\frac{T^n}{n}\right)
  =\prod_{i=0}^rP_i(T)^{(-1)^{i+1}}
  \qquad\text{in }E[[T]],
\]
where a negative exponent denotes the inverse in $E[[T]]$, available because
$P_i(0)=1$.  Equivalently, clearing the inverses,
\[
  \operatorname{Exp}\left(\sum_{n\geq1}N_n\frac{T^n}{n}\right)
  \cdot\prod_{\substack{0\leq i\leq r\\ i\text{ even}}}P_i(T)
  =\prod_{\substack{0\leq i\leq r\\ i\text{ odd}}}P_i(T),
\]
which is the shape of the factorization required by
\texttt{IsWeilFactorization}.  In particular the left-hand side is the Taylor
expansion at the origin of a rational function in $E(T)$.
\end{proposition}
\begin{proof}
By \Cref{lem:zla-counting-log} the exponent may be rewritten as the finite
alternating sum $\sum_i(-1)^i\sum_{n\geq1}\Tr(F_i^n\mid V_i)T^n/n$.  The formal
exponential turns a finite sum into a finite product, so the left-hand side is
$\prod_i\operatorname{Exp}\left(\sum_n\Tr(F_i^n)T^n/n\right)^{(-1)^i}$, and
each factor is $P_i(T)^{-1}$ by \Cref{lem:zla-trace-det}.
\end{proof}

\section{Rational germs and Hankel descent}\label{sec:zla-hankel}

\begin{definition}[Normalized rational germ]
\label{def:zla-rational-germ}
Let $K$ be a field.  A \emph{normalized polynomial fraction} is a pair
$(A,B)\in K[T]^2$ with $B(0)=1$.  Two pairs are equivalent when
\[
  (A,B)\sim(A',B')
  \quad\Longleftrightarrow\quad
  AB'=A'B.
\]
The resulting equivalence class is a \emph{normalized rational germ at zero}.
Its rational-function value is $A/B\in K(T)$ and its Taylor series is
\[
  \tau_K[A,B]=A(T)B(T)^{-1}\in K[[T]],
\]
the inverse being taken in $K[[T]]$.  Write $K(T)_0\subseteq K(T)$ for the
subring of rational functions regular at $T=0$.
\end{definition}

The Lean development does not package normalized rational germs as a type:
it works throughout with the predicate $\exists A,B\in K[T]$ with $B(0)=1$
and $fB=A$ on a power series $f$, so this definition has no Lean counterpart
of its own and the statements below are formalised in that unpackaged form.

\begin{lemma}[Faithfulness of the Taylor map]
\label{lem:zla-taylor-faithful}
\uses{def:zla-rational-germ}
\lean{PowerSeries.mul_eq_mul_of_mul_coe_eq,
      PowerSeries.eq_of_mul_eq_mul}
\leanok
If two normalized rational germs have the same Taylor series, then they are
equal.  Equivalently, Taylor expansion gives an injective map
\[
  K(T)_0\hookrightarrow K[[T]].
\]
\end{lemma}
\begin{proof}
\leanok
A denominator with constant coefficient $1$ is invertible in $K[[T]]$, so the
equality $AB^{-1}=A'B'^{-1}$ of power series may be cross-multiplied to
$AB'=A'B$ in $K[[T]]$; both sides are polynomials, and $K[T]\to K[[T]]$ is
injective, so the equality holds in $K[T]$, which is the defining relation of
the equivalence.
\end{proof}

\begin{definition}[Shift span and Hankel matrix]
\label{def:zla-shift-span}
For a sequence $a=(a_n)_{n\geq0}\in K^{\mathbf N}$, define the left shift $Sa$
by $(Sa)_n=a_{n+1}$, and put
\[
  \operatorname{Sh}_K(a)=\Span_K\{S^ma:m\geq0\}\subseteq K^{\mathbf N}.
\]
The associated Hankel matrix is $H(a)=(a_{r+s})_{r,s\geq0}$; its columns are
the shifts of $a$.  It has rank at most $h$ if every $(h+1)\times(h+1)$ minor
$\det(a_{r_x+s_y})_{0\leq x,y\leq h}$ vanishes, for all choices of nonnegative
indices $r_x,s_y$.
\end{definition}

Only the shift span is formalised; the Hankel matrix $H(a)$ and its rank
condition have no Lean counterpart, the shift span playing their role
throughout the Lean proofs.

\begin{definition}[Eventual recurrence]
\label{def:zla-eventual-recurrence}
An \emph{eventual recurrence of order $d$} for a sequence
$a\in K^{\mathbf N}$ is a tuple $(c_1,\ldots,c_d;N)$ in
$K^d\times\mathbf N$ satisfying
\[
  a_n+c_1a_{n-1}+\cdots+c_da_{n-d}=0
  \qquad(n\geq\max\{N,d\}).
\]
\end{definition}

There is no Lean definition of an eventual recurrence: it enters as the
hypothesis of \texttt{finiteDimensional\_shiftSpan\_of\_recurrence}.

\begin{lemma}[A rational germ supplies a recurrence]
\label{lem:zla-germ-recurrence}
\uses{def:zla-rational-germ, def:zla-eventual-recurrence}
\lean{PowerSeries.sum_coeff_mul_coeff_eq_zero}
\leanok
Let $f(T)=\sum_{n\geq0}a_nT^n$ be the Taylor series of the normalized rational
germ $A(T)/B(T)$, where $B(T)=1+c_1T+\cdots+c_dT^d$.  Then
\[
  a_n+c_1a_{n-1}+\cdots+c_da_{n-d}=0
\]
whenever $n>\deg A$ and $n\geq d$.  In particular $(c_1,\ldots,c_d;\deg A+1)$
is an eventual recurrence of order $d$ for $a$.
\end{lemma}
\begin{proof}
\leanok
Read off the coefficient of $T^n$ in the identity $B(T)f(T)=A(T)$: the
left-hand side contributes $\sum_{k=0}^{d}c_ka_{n-k}$ with $c_0=1$, and the
right-hand side contributes $0$ once $n>\deg A$.
\end{proof}

\begin{lemma}[A recurrence gives finite shift span]
\label{lem:zla-recurrence-shift-span}
\uses{def:zla-eventual-recurrence, def:zla-shift-span}
\lean{finiteDimensional_shiftSpan_of_recurrence}
\leanok
If $a\in K^{\mathbf N}$ satisfies an eventual recurrence, then
$\operatorname{Sh}_K(a)$ is finite-dimensional.
\end{lemma}
\begin{proof}
\leanok
Let $(c_1,\ldots,c_d;N)$ be the recurrence and put
$M=\max\{N,d\}$.  Beyond the threshold the recurrence expresses $S^ma$, for
$m\geq M$, as a $K$-linear combination of $S^{m-1}a,\ldots,S^{m-d}a$ up to a
sequence supported in the finitely many coordinates below the threshold; hence
$\operatorname{Sh}_K(a)$ is contained in the span of
$a,Sa,\ldots,S^{M-1}a$ together with the finitely many coordinate vectors
supported before the recurrence threshold, and by induction on $m$ every shift
lies in that finite-dimensional space.
\end{proof}

\begin{lemma}[Finite coordinates detect independent sequences]
\label{lem:zla-finite-coordinates}
\lean{linearIndependent_iff_exists_det_ne_zero,
      LinearIndependent.exists_det_ne_zero}
\leanok
Let $v_1,\ldots,v_r\in K^{\mathbf N}$ be linearly independent.  There exist
indices $n_1,\ldots,n_r\in\mathbf N$ such that
\[
  \det\left(v_j(n_i)\right)_{1\leq i,j\leq r}\neq0.
\]
\end{lemma}
\begin{proof}
\leanok
Induct on $r$ by finite Gaussian elimination: given indices that work for
$v_1,\ldots,v_{r-1}$, subtract from $v_r$ the unique combination of
$v_1,\ldots,v_{r-1}$ agreeing with it at those indices; the residual sequence
is nonzero by independence, so it does not vanish at some further coordinate
$n_r$, and expanding the determinant along the last row gives a nonzero
value.
\end{proof}

\begin{proposition}[Finite shift span supplies a rational germ]
\label{prop:zla-shift-span-germ}
\uses{def:zla-shift-span, def:zla-rational-germ}
Assume $W=\operatorname{Sh}_K(a)$ is finite-dimensional of dimension $r$.  The
shift restricts to an endomorphism $S_W:W\to W$.  Write
\[
  \det(U-S_W)=U^r+c_{r-1}U^{r-1}+\cdots+c_0
\]
and set $B_a(T)=1+c_{r-1}T+\cdots+c_0T^r$.  If
$f_a(T)=\sum_{n\geq0}a_nT^n$, then $B_a(T)f_a(T)$ is a polynomial $A_a(T)$.
If $r=0$ then $a=0$ and $A_a=0$; if $r>0$ then $\deg A_a<r$.  In both cases
$f_a(T)=\tau_K[A_a,B_a]$.
\end{proposition}
\begin{proof}
Cayley--Hamilton applied to $S_W$ gives
$S^ra+c_{r-1}S^{r-1}a+\cdots+c_0a=0$ in $W$, that is
$a_{n+r}+c_{r-1}a_{n+r-1}+\cdots+c_0a_n=0$ for every $n\geq0$.  This is
exactly the vanishing of the coefficient of $T^{n+r}$ in $B_a(T)f_a(T)$, so
that product is a polynomial of degree less than $r$.  Since $B_a(0)=1$, the
pair $(A_a,B_a)$ is a normalized polynomial fraction with Taylor series
$f_a$.
\end{proof}

The Lean statement asserts the existence of the normalized fraction
$(A_a,B_a)$ only; the degree bound $\deg A_a<r$ is visible in its proof but
is not part of the statement, and nothing below uses it.

\begin{lemma}[Finite determinants reflect scalar extension]
\label{lem:zla-det-scalar-extension}
\uses{lem:zla-finite-coordinates}
Let $K\hookrightarrow L$ be an injective field homomorphism.  For every finite
square matrix $M$ over $K$,
\[
  \det(M)=0
  \quad\Longleftrightarrow\quad
  \det(M\otimes_KL)=0.
\]
Consequently a finite family of sequences in $K^{\mathbf N}$ that is linearly
independent over $K$ remains linearly independent over $L$, and every finite
bound on the dimension of the shift span over $L$ descends to a bound over
$K$.
\end{lemma}
\begin{proof}
The determinant is a finite polynomial expression in the matrix entries and is
compatible with the ring homomorphism $K\to L$, which is injective, so it
reflects the vanishing of that expression.  The consequence follows from
\Cref{lem:zla-finite-coordinates}: a $K$-independent family has a nonvanishing
finite coordinate determinant, which stays nonzero over $L$.
\end{proof}

The determinant criterion itself is not isolated as a Lean declaration; the
two Lean statements are the two consequences, the first as an equivalence.

\begin{theorem}[Hankel rationality descent]
\label{thm:zla-hankel-descent}
\uses{def:zla-rational-germ, def:zla-shift-span, lem:zla-germ-recurrence,
lem:zla-recurrence-shift-span, lem:zla-finite-coordinates,
lem:zla-det-scalar-extension, prop:zla-shift-span-germ}
\lean{PowerSeries.exists_normalized_fraction_map_iff}
\leanok
For every field extension $K\subseteq L$,
\[
  K[[T]]\cap L(T)_0=K(T)_0
\]
inside $L[[T]]$.  Explicitly, if $f\in K[[T]]$ becomes the Taylor series of a
normalized rational germ over $L$, then its coefficient sequence has finite
shift span over $K$, and the characteristic polynomial of the shift on that
span constructs a normalized rational germ over $K$ whose Taylor series is
$f$.
\end{theorem}
\begin{proof}
\leanok
Let $a$ be the coefficient sequence of $f$, which lies in $K^{\mathbf N}$.  A
normalized rational germ over $L$ with Taylor series $f$ supplies an eventual
recurrence for $a$ over $L$ (\Cref{lem:zla-germ-recurrence}), hence a finite
bound on $\dim_L\operatorname{Sh}_L(a)$
(\Cref{lem:zla-recurrence-shift-span}).  Any $K$-independent family of shifts
of $a$ has a nonvanishing finite coordinate determinant
(\Cref{lem:zla-finite-coordinates}) and therefore stays independent over $L$
(\Cref{lem:zla-det-scalar-extension}); so $\dim_K\operatorname{Sh}_K(a)$ is
bounded by the same number.  Now \Cref{prop:zla-shift-span-germ} over $K$
produces $A_a,B_a\in K[T]$ with $B_a(0)=1$ and
$f=\tau_K[A_a,B_a]$.  The reverse inclusion is trivial.
\end{proof}

\begin{corollary}[Rationality over the counting field]
\label{cor:zla-rationality-counting-field}
\uses{thm:zla-hankel-descent, def:zla-rational-germ}
\lean{PowerSeries.exists_normalized_fraction_of_map_rat}
\leanok
If $Z(T)\in\mathbf Q[[T]]$ becomes the Taylor series of a normalized rational
germ after coefficientwise extension to a characteristic-zero field $E$, then
$Z(T)\in\mathbf Q(T)_0$: there are $A,B\in\mathbf Q[T]$ with $B(0)=1$ and
$B(T)Z(T)=A(T)$.
\end{corollary}
\begin{proof}
\leanok
Apply \Cref{thm:zla-hankel-descent} to the extension $\mathbf Q\subseteq E$,
which is an extension of fields because $E$ has characteristic zero.
\end{proof}

\section{Integral Fatou and normalized Gauss descent}\label{sec:zla-fatou}

\begin{definition}[Integral shift lattice]
\label{def:zla-integral-shift-lattice}
\uses{def:zla-shift-span}
\lean{intSeqShift,
      intShiftSpan}
\leanok
Let $a\in\mathbf Z^{\mathbf N}$ have finite-dimensional rational shift span
$W=\operatorname{Sh}_{\mathbf Q}(a)$.  Define
\[
  M(a)=\sum_{m\geq0}\mathbf Z\,S^ma\subseteq W.
\]
This is a shift-stable torsion-free abelian group with
$M(a)\otimes_{\mathbf Z}\mathbf Q=W$.
\end{definition}

\begin{lemma}[Finite-coordinate lattice]
\label{lem:zla-finite-coordinate-lattice}
\uses{def:zla-integral-shift-lattice, lem:zla-finite-coordinates}
\lean{Submodule.exists_finset_eval_injective,
      finite_intShiftSpan,
      free_intShiftSpan}
\leanok
With the notation of \Cref{def:zla-integral-shift-lattice}, there is a finite
set $J\subset\mathbf N$ such that coordinate evaluation
\[
  \operatorname{ev}_J:W\longrightarrow\mathbf Q^J
\]
is injective.  Its restriction maps $M(a)$ injectively into $\mathbf Z^J$.
Hence $M(a)$ is a finite free abelian group.
\end{lemma}
\begin{proof}
\leanok
Choose a basis of $W$ consisting of shifts of $a$ and let $J$ be a set of
coordinates on which the corresponding determinant does not vanish
(\Cref{lem:zla-finite-coordinates}); then $\operatorname{ev}_J$ is injective on
$W$ for dimension reasons.  Every element of $M(a)$ is an integral linear
combination of integral sequences, hence integral, so
$\operatorname{ev}_J(M(a))\subseteq\mathbf Z^J$, and a subgroup of a finite
free abelian group is finite free.
\end{proof}

\begin{lemma}[Integral Cayley--Hamilton recurrence]
\label{lem:zla-integral-cayley-hamilton}
\uses{def:zla-integral-shift-lattice, lem:zla-finite-coordinate-lattice}
\lean{exists_monic_sum_coeff_mul_eq_zero}
\leanok
With the notation of \Cref{def:zla-integral-shift-lattice}, the shift
endomorphism $S:M(a)\to M(a)$ has a monic characteristic polynomial
\[
  C(U)=U^r+c_{r-1}U^{r-1}+\cdots+c_0\in\mathbf Z[U],
\]
where $r$ is the rank of $M(a)$, and for every $n\geq0$,
\[
  a_{n+r}+c_{r-1}a_{n+r-1}+\cdots+c_0a_n=0.
\]
\end{lemma}
\begin{proof}
\leanok
By \Cref{lem:zla-finite-coordinate-lattice} the group $M(a)$ is finite free, and
it is stable under the shift, so $S$ has an integral matrix in any basis; its
characteristic polynomial is monic with integer coefficients.
Cayley--Hamilton over $\mathbf Z$ gives
$S^ra+c_{r-1}S^{r-1}a+\cdots+c_0a=0$ in $M(a)$, and reading this equality of
sequences in coordinate $n$ is the displayed recurrence.
\end{proof}

\begin{proposition}[Integral rational presentation]
\label{prop:zla-integral-presentation}
\uses{lem:zla-integral-cayley-hamilton, thm:zla-hankel-descent,
def:zla-rational-germ}
\lean{PowerSeries.exists_int_normalized_fraction'}
\leanok
Let $f(T)=\sum_{n\geq0}a_nT^n\in\mathbf Z[[T]]$ be the Taylor series of a
normalized rational germ over $\mathbf Q$.  Then there exist
$A_0,B_0\in\mathbf Z[T]$ with $B_0(0)=1$ and $A_0(0)=a_0$ such that
\[
  B_0(T)f(T)=A_0(T)
  \qquad\text{in }\mathbf Z[[T]].
\]
One may take $B_0(T)=1+c_{r-1}T+\cdots+c_0T^r$ built from the characteristic
polynomial of the shift on the integral shift lattice.
\end{proposition}
\begin{proof}
\leanok
Rationality over $\mathbf Q$ makes $\operatorname{Sh}_{\mathbf Q}(a)$
finite-dimensional, so the integral shift lattice is defined and
\Cref{lem:zla-integral-cayley-hamilton} applies.  Multiplying $f$ by the
integral polynomial $B_0$ built from the recurrence kills every coefficient in
degree at least $r$, so the product is a polynomial with integer coefficients;
its constant coefficient is $a_0$ because $B_0(0)=1$.
\end{proof}

\begin{lemma}[Normalized Gauss factor lemma]
\label{lem:zla-normalized-gauss}
\lean{Polynomial.exists_int_of_coeff_zero_eq_one_of_mul}
\leanok
Let $A\in\mathbf Z[T]$ satisfy $A(0)=1$.  If
\[
  A(T)=P(T)Q(T),\qquad P,Q\in\mathbf Q[T],\qquad P(0)=Q(0)=1,
\]
then $P,Q\in\mathbf Z[T]$.
\end{lemma}
\begin{proof}
\leanok
Reversal turns a polynomial with constant coefficient $1$ and degree $b$ into
the monic polynomial $U^bA(U^{-1})$, and it is multiplicative for the
degrees involved here, so the factorization becomes a factorization of the
monic integral polynomial $U^{\deg A}A(U^{-1})$ into the monic rational
polynomials $U^{\deg P}P(U^{-1})$ and $U^{\deg Q}Q(U^{-1})$.  Gauss's lemma --
primitive integral normalization together with multiplicativity of the content
-- forces both monic rational factors to be integral, and reversing back gives
$P,Q\in\mathbf Z[T]$.
\end{proof}

\begin{proposition}[Normalized gcd reduction over the integers]
\label{prop:zla-normalized-gcd-reduction}
\uses{lem:zla-normalized-gauss}
\lean{Polynomial.exists_reduced_int_fraction}
\leanok
Let $A_0,B_0\in\mathbf Z[T]$ satisfy $A_0(0)=B_0(0)=1$, and let
$G\in\mathbf Q[T]$ be the unique generator of the ideal
$(A_0,B_0)\subseteq\mathbf Q[T]$ with $G(0)=1$.  Define
$A,B\in\mathbf Q[T]$ by $A_0=GA$ and $B_0=GB$.  Then
\[
  G,A,B\in\mathbf Z[T],
  \qquad
  G(0)=A(0)=B(0)=1,
\]
the pair $(A,B)$ is coprime in $\mathbf Q[T]$, that is, every common divisor
of $A$ and $B$ is a unit, and
\[
  \frac{A_0}{B_0}=\frac{A}{B}
  \qquad\text{in }\mathbf Q(T).
\]
\end{proposition}
\begin{proof}
\leanok
The ring $\mathbf Q[T]$ is a principal ideal domain, so the ideal
$(A_0,B_0)$ has a generator, unique up to a scalar; since $A_0(0)=1$ the
generator does not vanish at $0$ and may be normalized by $G(0)=1$.  Then
$A(0)=B(0)=1$ follows from $A_0(0)=B_0(0)=1$, and
\Cref{lem:zla-normalized-gauss} applied to $A_0=GA$ and to $B_0=GB$ gives
$G,A,B\in\mathbf Z[T]$.  A common divisor of $A$ and $B$ would multiply $G$ to
a larger common divisor of $A_0$ and $B_0$, contradicting that $G$ generates
their ideal.  The fraction is unchanged because $G\neq0$.
\end{proof}

The Lean statement asserts the existence of the reduced pair $(A,B)$ with
the stated properties; it does not name the common factor $G$ or record its
integrality separately.

\begin{lemma}[Uniqueness of normalized reduced fractions]
\label{lem:zla-normalized-fraction-unique}
\lean{Polynomial.eq_of_coeff_zero_eq_one_of_isCoprime}
\leanok
Let $K$ be a field.  Suppose $A,B,O,E\in K[T]$ have constant coefficient $1$,
both pairs $(A,B)$ and $(O,E)$ are coprime, and $OB=EA$.  Then
\[
  O=A,\qquad E=B.
\]
\end{lemma}
\begin{proof}
\leanok
From $OB=EA$ and coprimality of $(A,B)$, Euclid's lemma gives $A\mid O$ and
$B\mid E$; symmetrically coprimality of $(O,E)$ gives $O\mid A$ and
$E\mid B$.  So $O$ and $A$ agree up to a unit of $K[T]$, that is a nonzero
scalar, and likewise for $E$ and $B$; the normalization by constant
coefficient $1$ makes both scalars equal to $1$.
\end{proof}

\begin{theorem}[Integral Fatou theorem]
\label{thm:zla-integral-fatou}
\uses{prop:zla-integral-presentation, prop:zla-normalized-gcd-reduction,
lem:zla-normalized-fraction-unique, def:zla-rational-germ}
\lean{PowerSeries.exists_unique_reduced_int_fraction}
\leanok
Let $f(T)\in1+T\mathbf Z[[T]]$ be the Taylor series of a normalized rational
germ over $\mathbf Q$.  There is a unique pair of polynomials
$A,B\in\mathbf Z[T]$, coprime in $\mathbf Q[T]$, satisfying
\[
  A(0)=B(0)=1,
  \qquad
  f(T)=\frac{A(T)}{B(T)}.
\]
This is the \emph{normalized reduced integral presentation} of $f$.
\end{theorem}
\begin{proof}
\leanok
\Cref{prop:zla-integral-presentation} produces $A_0,B_0\in\mathbf Z[T]$ with
$B_0f=A_0$; here $A_0(0)=f(0)=1=B_0(0)$.  Reducing this fraction by
\Cref{prop:zla-normalized-gcd-reduction} gives coprime integral $A,B$ with
constant coefficient $1$ and $A/B=A_0/B_0=f$.  If $(A',B')$ is a second pair
with the same properties, then $A'/B'=f=A/B$ and both denominators are nonzero,
so cross multiplication gives $A'B=B'A$; both pairs are coprime and all four
polynomials have constant coefficient $1$, so
\Cref{lem:zla-normalized-fraction-unique} over $\mathbf Q$ yields $A'=A$ and
$B'=B$.
\end{proof}

\begin{definition}[Reversal and reciprocal roots]
\label{def:zla-reversal}
\lean{Polynomial.reverse}
\leanok
Let $P(T)=1+c_1T+\cdots+c_bT^b\in\mathbf Z[T]$ have degree $b$.  Its
\emph{reversal} is
\[
  P^{\mathrm{rev}}(U)=U^bP(U^{-1})=U^b+c_1U^{b-1}+\cdots+c_b\in\mathbf Z[U].
\]
For a field $\Omega$ containing $\mathbf Q$, a scalar
$\alpha\in\Omega^\times$ is a \emph{reciprocal root} of $P$ if
$P(\alpha^{-1})=0$.  The same definitions apply verbatim to a polynomial with
constant coefficient $1$ over any field containing $\mathbf Q$.
\end{definition}

\begin{lemma}[Reciprocal roots and reversal]
\label{lem:zla-reversal-roots}
\uses{def:zla-reversal}
\lean{Polynomial.reverse_leadingCoeff,
      Polynomial.eval₂_reverse_eq_zero_iff}
\leanok
For $P$ as in \Cref{def:zla-reversal}, the polynomial $P^{\mathrm{rev}}$ is
monic, and for every field $\Omega$ containing $\mathbf Q$ and every
$\alpha\in\Omega^\times$,
\[
  P(\alpha^{-1})=0
  \quad\Longleftrightarrow\quad
  P^{\mathrm{rev}}(\alpha)=0.
\]
\end{lemma}
\begin{proof}
\leanok
Monicity is the coefficient $c_0=1$ of $P$ read in top degree after reversal.
The equivalence is the defining identity
$P^{\mathrm{rev}}(\alpha)=\alpha^bP(\alpha^{-1})$ together with
$\alpha\neq0$.
\end{proof}

\section{Weights and recovery of individual factors}\label{sec:zla-weights}

Fix an integer $q>1$ and let $\Omega$ be an algebraically closed field of
characteristic zero; the case of interest is an abstract algebraic closure of
$\Qell$.  Where the sketch article measures the size of a scalar through a
\emph{universal complex-realization field} and its complex realizations, that
field and its realizations are replaced here by $\Qbar\subseteq\mathbf C$ and
the conjugates of the scalar, because
\texttt{AlgebraicGeometry.WeilRiemannHypothesis} of
\texttt{Comparator/Weil/Definitions.lean} speaks about complex roots.

In the Lean formalisation of this section the field $\Omega$ is $\mathbf C$ and
the weights are natural numbers; the statements tagged below are those
specialisations.

\begin{definition}[Universal purity of a scalar]
\label{def:zla-universal-purity}
A nonzero algebraic number $\alpha\in\Qbar\subseteq\mathbf C$ is
\emph{universally pure of weight $w\in\mathbf R$} if every complex root of the
minimal polynomial of $\alpha$ over $\mathbf Q$ has absolute value $q^{w/2}$;
equivalently, $|\tau(\alpha)|=q^{w/2}$ for every embedding
$\tau:\mathbf Q(\alpha)\hookrightarrow\mathbf C$, these embeddings being the
complex roots of that minimal polynomial.  The condition is read off the
minimal polynomial alone, so it applies verbatim to a nonzero element of any
field of characteristic zero that is algebraic over $\mathbf Q$.
\end{definition}

The Lean predicate is stated for the inverse scalar:
\texttt{IsUniversallyPureOfWeight q w z} says that every complex root of the
minimal polynomial of $z$ has absolute value $q^{-w/2}$, that is, that
$z^{-1}$ is universally pure of weight $w$ in the sense above.

\begin{lemma}[Uniqueness of a scalar weight]
\label{lem:zla-scalar-weight-unique}
\uses{def:zla-universal-purity}
\lean{IsUniversallyPureOfWeight.weight_unique}
\leanok
If a nonzero algebraic number is universally pure of weight $w$ and of weight
$w'$, then $w=w'$.
\end{lemma}
\begin{proof}
\leanok
The scalar is itself a complex root of its own minimal polynomial, so its
absolute value is both $q^{w/2}$ and $q^{w'/2}$; the real exponential with
base $q>1$ is injective, so $w=w'$.
\end{proof}

\begin{definition}[Weight stratum of a rational polynomial]
\label{def:zla-weight-stratum}
\lean{Polynomial.weightFactor,
      Polynomial.coweightFactor}
\leanok
Let $R\in\mathbf Q[T]$ satisfy $R(0)=1$.  Since $R(0)\neq0$, no irreducible
factor of $R$ is a scalar multiple of $T$, so $R$ has a factorization
\[
  R=\prod_{j}F_j^{m_j},
  \qquad
  F_j\in\mathbf Q[T]\text{ irreducible},\quad F_j(0)=1,
\]
with the $F_j$ pairwise distinct and $m_j\geq1$, and this normalized
irreducible factorization is unique.  For a real number $w$, the
\emph{weight-$w$ stratum} of $R$ is
\[
  R_w(T)=\prod_{j\in J_w}F_j(T)^{m_j},
  \qquad
  J_w=\{j:|z|=q^{-w/2}\text{ for every complex root }z\text{ of }F_j\},
\]
the product of those normalized irreducible rational factors of $R$ all of
whose complex roots have absolute value $q^{-w/2}$, taken with the
multiplicities they have in $R$.
\end{definition}

\begin{proposition}[Rationality of a weight stratum]
\label{prop:zla-weight-stratum-rational}
\uses{def:zla-weight-stratum, def:zla-universal-purity, def:zla-reversal,
lem:zla-reversal-roots}
Let $R\in\mathbf Q[T]$ satisfy $R(0)=1$ and let $w$ be a real number.  Then
$R_w\in\mathbf Q[T]$ divides $R$, both $R_w$ and the complementary factor
$R/R_w\in\mathbf Q[T]$ have constant coefficient $1$, and for every field $K$
of characteristic zero the reciprocal roots of $R_w$ in $K$, counted with
multiplicity, are exactly the reciprocal roots of $R$ in $K$ that are
universally pure of weight $w$.  In particular, for $K=\mathbf C$, the
complex roots of $R_w$ are exactly the complex roots of $R$ whose inverses are
universally pure of weight $w$.
\end{proposition}
\begin{proof}
Rationality is by construction: $R_w$ is a subproduct of the normalized
irreducible factorization of $R$ over $\mathbf Q$, so it lies in
$\mathbf Q[T]$, divides $R$, and its complementary factor is the product of
the remaining $F_j^{m_j}$, hence also lies in $\mathbf Q[T]$; both constant
coefficients are $1$ because every $F_j$ has constant coefficient $1$.

For the description of the roots it suffices to see that the purity condition
is constant on each irreducible factor.  Let $F$ be one of the $F_j$ and let
$\alpha\in K^\times$ be a reciprocal root of $F$, that is $F(\alpha^{-1})=0$.
By \Cref{lem:zla-reversal-roots} this says $F^{\mathrm{rev}}(\alpha)=0$, and
$F^{\mathrm{rev}}$ is monic; it is irreducible in $\mathbf Q[U]$ because $F$
is irreducible with $F(0)\neq0$.  So $\alpha$ is algebraic over $\mathbf Q$
with minimal polynomial $F^{\mathrm{rev}}$, whose complex roots are exactly
the inverses of the complex roots of $F$.  Hence $\alpha$ is universally pure
of weight $w$ if and only if $|z|=q^{-w/2}$ for every complex root $z$ of $F$,
a condition on $F$ alone.  So the reciprocal roots of $R$ in $K$ that are
universally pure of weight $w$ are exactly those contributed by the factors
$F_j$ with $j\in J_w$, with the multiplicities $m_j$ they have in $R$, and
these are the reciprocal roots of $R_w$.
\end{proof}

The Lean description of the roots is the case $K=\mathbf C$, that is the
last sentence of the statement; the general characteristic-zero field $K$ is
not formalised.

Purity under a single embedding would not suffice here: the reciprocal roots
of $R$ of one prescribed absolute value need not form a conjugation-stable
family, and the factor they cut out need not be rational, which is why
\Cref{def:zla-universal-purity} tests every conjugate.

\begin{definition}[Pure polynomial]
\label{def:zla-pure-polynomial}
\uses{def:zla-reversal, def:zla-universal-purity}
\lean{Polynomial.IsUniversallyPureOfWeight}
\leanok
A polynomial $P(T)\in\Omega[T]$ with $P(0)=1$ is \emph{pure of weight $w$} if
all of its reciprocal roots, counted with multiplicity, are nonzero, algebraic
over $\mathbf Q$, and universally pure of weight $w$ in the sense of
\Cref{def:zla-universal-purity}.
\end{definition}

For the rest of this section fix $r\geq0$ and polynomials
$P_i(T)\in\Omega[T]$, $0\leq i\leq r$, with $P_i(0)=1$, and put
\[
  O(T)=\prod_{\substack{0\leq i\leq r\\ i\text{ odd}}}P_i(T),
  \qquad
  E(T)=\prod_{\substack{0\leq i\leq r\\ i\text{ even}}}P_i(T).
\]

\begin{lemma}[Different weights have disjoint reciprocal roots]
\label{lem:zla-disjoint-weights}
\uses{def:zla-pure-polynomial, lem:zla-scalar-weight-unique}
\lean{Polynomial.IsPureOfWeight.disjoint_roots}
\leanok
Assume $P_i$ is pure of weight $i$ for every $i$.  If $i\neq j$, then $P_i$
and $P_j$ have no common reciprocal root in $\Omega$.
\end{lemma}
\begin{proof}
\leanok
A common reciprocal root would be universally pure of weight $i$ and of weight
$j$, so $i=j$ by \Cref{lem:zla-scalar-weight-unique}.
\end{proof}

The Lean statement, and likewise that of
\Cref{prop:zla-parity-identification}, assumes only that the complex roots
of each $P_i$ have absolute value $q^{-i/2}$, which is weaker than the
universal purity of \Cref{def:zla-pure-polynomial}; so the Lean statement
proves more.

\begin{proposition}[Identification of the two parity factors]
\label{prop:zla-parity-identification}
\uses{def:zla-pure-polynomial, lem:zla-disjoint-weights,
lem:zla-normalized-fraction-unique}
\lean{Polynomial.prod_odd_eq_map_and_prod_even_eq_map}
\leanok
Assume every $P_i$ is pure of weight $i$.  Let $A,B\in\mathbf Q[T]$ be
coprime with $A(0)=B(0)=1$, and suppose that
\[
  \frac{O(T)}{E(T)}=\frac{A(T)}{B(T)}
  \qquad\text{in }\Omega(T).
\]
Then, after coefficient extension to $\Omega$,
\[
  O(T)=A(T),\qquad E(T)=B(T).
\]
Thus no reciprocal root cancels between odd and even cohomological degrees.
\end{proposition}
\begin{proof}
\leanok
By \Cref{lem:zla-disjoint-weights} the polynomials $O$ and $E$ have no common
root in the algebraically closed field $\Omega$, hence are coprime in
$\Omega[T]$; the pair $(A,B)$ stays coprime after extension of scalars since a
Bézout identity extends.  Both denominators are nonzero because their constant
coefficients are $1$, so the equality of fractions may be cross-multiplied to
$OB=EA$ in $\Omega[T]$, and all four polynomials have constant coefficient
$1$.  \Cref{lem:zla-normalized-fraction-unique} over $\Omega$ concludes.
\end{proof}

\begin{definition}[Complete weight factor]
\label{def:zla-complete-weight-factor}
\uses{def:zla-weight-stratum, prop:zla-weight-stratum-rational}
Continue with the hypotheses above and assume $A,B\in\mathbf Z[T]$.  For
$0\leq i\leq r$ put
\[
  R_i(T)=
  \begin{cases}
    A(T),&i\text{ odd},\\
    B(T),&i\text{ even},
  \end{cases}
\]
and define the \emph{complete weight-$i$ factor} $P_i^{\mathrm{rec}}(T)$ of
the reduced fraction $A/B$ to be the weight-$i$ stratum of $R_i$ in the sense
of \Cref{def:zla-weight-stratum}.  By
\Cref{prop:zla-weight-stratum-rational} it lies in $\mathbf Q[T]$, satisfies
$P_i^{\mathrm{rec}}(0)=1$, and its reciprocal roots are exactly the
reciprocal roots of $R_i$ that are universally pure of weight $i$, with the
multiplicities they have in $R_i$.
\end{definition}

\begin{proposition}[Unique recovery over the coefficient field]
\label{prop:zla-unique-recovery}
\uses{def:zla-complete-weight-factor, prop:zla-parity-identification,
def:zla-pure-polynomial, prop:zla-weight-stratum-rational}
For every $i$,
\[
  P_i(T)=P_i^{\mathrm{rec}}(T)
  \qquad\text{in }\Omega[T].
\]
In particular $P_i$ is the scalar extension of a unique polynomial in
$\mathbf Q[T]$, determined by the reduced fraction $A/B$, by $q$, and by the
integer $i$.
\end{proposition}
\begin{proof}
By \Cref{prop:zla-parity-identification} the product of the $P_j$ over the
$j$ of the same parity as $i$ is the scalar extension of $R_i$, so the
reciprocal-root multiset of $R_i$ is the union, with multiplicities, of those
of the $P_j$ with $j$ of that parity.  Purity and \Cref{lem:zla-disjoint-weights} show that the roots coming
from $P_i$ are exactly the reciprocal roots of $R_i$ that are universally pure
of weight $i$, and by \Cref{prop:zla-weight-stratum-rational} these are
exactly the reciprocal roots of the weight-$i$ stratum of $R_i$, with the same
multiplicities; hence $P_i$ and $P_i^{\mathrm{rec}}$ have the same
reciprocal-root multiset.  Two normalized polynomials with the same
reciprocal-root multiset and constant coefficient $1$ are equal.
\end{proof}

\begin{proposition}[Integrality of an individual factor]
\label{prop:zla-factor-integrality}
\uses{def:zla-complete-weight-factor, prop:zla-weight-stratum-rational,
lem:zla-normalized-gauss}
The recovered polynomial $P_i^{\mathrm{rec}}(T)$ belongs to $\mathbf Z[T]$.
\end{proposition}
\begin{proof}
By \Cref{prop:zla-weight-stratum-rational} the polynomial
$P_i^{\mathrm{rec}}$ is a factor of $R_i$ in $\mathbf Q[T]$ whose
complementary factor $R_i/P_i^{\mathrm{rec}}$ is again rational, and both
have constant coefficient $1$.  Now $R_i$ is $A$ or $B$, hence lies in
$\mathbf Z[T]$ with $R_i(0)=1$, and \Cref{lem:zla-normalized-gauss} applied to
this factorization gives $P_i^{\mathrm{rec}}\in\mathbf Z[T]$.
\end{proof}

\section{Perfect similitude pairings}\label{sec:zla-similitude}

\begin{definition}[Perfect similitude pairing]
\label{def:zla-similitude-pairing}
Let $E$ be a field and let $V,W$ be finite-dimensional $E$-vector spaces.  A
bilinear pairing $B:V\times W\to E$ is \emph{perfect} if the adjoint map
\[
  \phi_B:W\longrightarrow V^\vee,
  \qquad
  \phi_B(w)(v)=B(v,w),
\]
is an isomorphism.  Given automorphisms $F_V\in\operatorname{Aut}_E(V)$,
$F_W\in\operatorname{Aut}_E(W)$ and a scalar $\mu\in E^\times$, the pairing is
a \emph{similitude pairing of multiplier $\mu$} if
\[
  B(F_Vv,F_Ww)=\mu B(v,w)
  \qquad(v\in V,\ w\in W).
\]
\end{definition}

Only perfectness of the pairing is a Lean definition; the similitude
identity is carried as the hypothesis \texttt{p (F v) (G w) = c * p v w} of
the statements below, not packaged as a definition.

\begin{lemma}[Adjoint conjugacy]
\label{lem:zla-adjoint-conjugacy}
\uses{def:zla-similitude-pairing}
\lean{LinearMap.dualMap_comp_toPerfPair_comp_of_similitude,
      LinearMap.charpoly_eq_smul_inverse_dualMap_of_similitude}
\leanok
Let $B:V\times W\to E$ be a perfect similitude pairing of multiplier $\mu$ for
the automorphisms $F_V,F_W$.  Then
\[
  \phi_BF_W=\mu\,(F_V^{-1})^\vee\phi_B,
\]
so $F_W$ is conjugate to $\mu\,(F_V^{-1})^\vee$.
\end{lemma}
\begin{proof}
\leanok
Evaluate both sides at $w\in W$ and then at $v\in V$: the left side gives
$B(v,F_Ww)$ and the right side gives $\mu B(F_V^{-1}v,w)$, and these agree by
the similitude identity applied to the pair $(F_V^{-1}v,w)$.  Since $\phi_B$
is an isomorphism, the displayed identity exhibits the conjugacy.
\end{proof}

\begin{proposition}[Reciprocal eigenvalue pairing]
\label{prop:zla-reciprocal-eigenvalue-pairing}
\uses{def:zla-similitude-pairing, lem:zla-adjoint-conjugacy}
Let $B:V\times W\to E$ be a perfect similitude pairing of multiplier $\mu$ for
the automorphisms $F_V,F_W$, let $\Omega$ be an algebraic closure of $E$, and
put $b=\dim_EV=\dim_EW$.  If $\alpha_1,\ldots,\alpha_b\in\Omega$ are the
eigenvalues of $F_V\otimes_E\Omega$, counted with algebraic multiplicity, then
the eigenvalues of $F_W\otimes_E\Omega$, counted with algebraic multiplicity,
are
\[
  \frac{\mu}{\alpha_1},\ldots,\frac{\mu}{\alpha_b}.
\]
No semisimplicity assumption is required.
\end{proposition}
\begin{proof}
By \Cref{lem:zla-adjoint-conjugacy} the endomorphism $F_W$ is conjugate to
$\mu\,(F_V^{-1})^\vee$, and conjugate endomorphisms have the same
characteristic polynomial.  Dualizing does not change the characteristic
polynomial, inverting an automorphism replaces its eigenvalue multiset by the
multiset of inverses, and multiplying by the scalar $\mu$ scales every
eigenvalue by $\mu$; all three operations are computed on characteristic
polynomials, so no semisimplicity is needed.
\end{proof}

The Lean statement is the conclusion over an algebraically closed base
field: the pairing is assumed to be one over $\Omega$ already, rather than
obtained from a pairing over $E$ by scalar extension.

\begin{theorem}[Degreewise reciprocal-polynomial identities]
\label{thm:zla-degreewise-reciprocal}
\uses{def:zla-similitude-pairing, prop:zla-reciprocal-eigenvalue-pairing,
def:zla-recip-charpoly}
\lean{LinearMap.C_det_mul_charpolyRev_of_similitude}
\leanok
Let $B:V\times W\to E$ be a perfect similitude pairing of multiplier $\mu$ for
the automorphisms $F_V,F_W$.  Put $b=\dim_EV=\dim_EW$ and
$D_V=\det(F_V)$, and set
\[
  P_V(T)=\det(1-TF_V\mid V),
  \qquad
  P_W(T)=\det(1-TF_W\mid W).
\]
Then, in $E[T,T^{-1}]$,
\[
  P_W(\mu^{-1}T^{-1})=(-1)^bT^{-b}D_V^{-1}P_V(T),
\]
or equivalently
\[
  P_W(T)=(-1)^b\mu^bD_V^{-1}T^bP_V(\mu^{-1}T^{-1}).
\]
\end{theorem}
\begin{proof}
\leanok
Over an algebraic closure, $P_V(T)=\prod_{j=1}^b(1-\alpha_jT)$ with the
$\alpha_j$ the eigenvalues of $F_V$, all nonzero, and
$D_V=\prod_j\alpha_j$.  By
\Cref{prop:zla-reciprocal-eigenvalue-pairing},
$P_W(T)=\prod_j(1-\mu\alpha_j^{-1}T)$.  Substituting $T\mapsto\mu^{-1}T^{-1}$
and clearing $\prod_j(-\alpha_j^{-1}T^{-1})$ from the product gives the first
identity; both sides have coefficients in $E$, so the identity descends from
$\Omega[T,T^{-1}]$ to $E[T,T^{-1}]$.  The second form follows by substituting
$T\mapsto\mu^{-1}T^{-1}$ once more.
\end{proof}
